\chapter{Budgeting and Financial Control}

Tracking spending tells me where money went. Budgeting tells me where it should go. This chapter builds the monthly budget that connects the Phase~1 spending data to the Phase~2 savings goals.

\section{Conceptual Definitions of Expenses}

\textbf{Fixed expenses} are regular, unchanging monthly costs. Mine are just 320,000~VND (internet: 250,000~VND; phone card: 70,000~VND) because my family covers apartment fees, electricity, and water. \textbf{Variable expenses} change based on daily choices: food, sports, shopping, and entertainment. \textbf{Planned expenses} are predictable future costs I can budget for in advance---racket restringing (approx.\ 440,000~VND/month) and shuttlecock purchases (approx.\ 300,000~VND/month). \textbf{Unplanned expenses} are genuine surprises: a clinic visit for a sprained ankle (300,000~VND) and a motorcycle turn signal replacement (300,000~VND) from the tracking period make the case for keeping a monthly emergency buffer.

\begin{table}[h!]
\caption{Classification of Personal Expenses by Type}
\label{tab:expense_classification}
\begin{tabularx}{\linewidth}{@{}p{3cm}p{1.8cm}Xr@{}}
\toprule
\textbf{Category} & \textbf{Type} & \textbf{Controllability} & \textbf{Current Monthly (VND)} \\
\midrule
Internet \& Phone Card  & Fixed    & None --- set by contract    & 320,000   \\
Food (groceries)        & Variable & Medium --- can plan meals   & $\approx$1,200,000 \\
Food (dining out)       & Variable & High --- discretionary      & $\approx$2,600,000 \\
Sports (court \& equipment) & Planned & Medium --- set schedule  & 759,000   \\
Shopping online         & Variable & High --- fully discretionary & 474,999  \\
Entertainment           & Variable & High --- fully discretionary & 145,000  \\
Unplanned (health, moto) & Unplanned & Low --- emergencies       & $\approx$300,000 \\
\midrule
\textbf{Total} & & & \textbf{5,798,999} \\
\bottomrule
\end{tabularx}
\end{table}

The classification reveals an important structural point: almost all of my spending falls into the variable or planned categories, which means almost all of it responds to conscious decisions. Fixed costs---the portion that cannot be changed without switching providers or canceling contracts---represent just 320,000~VND out of roughly 5.8 million in total spending, or 5.5\%. This is unusually low and is almost entirely a result of the family housing arrangement. It also means that virtually every line in my budget is improvable through behavior, not just circumstance. The only category that resists behavioral control is the unplanned one---which is precisely why the 300,000~VND monthly emergency buffer in the target budget exists: not to predict what will go wrong, but to ensure that when something does, it does not come out of next month's savings.



\section{Monthly Target Budget}

\begin{table}[H]
\caption{Monthly Target Budget Structure}
\label{tab:target_budget}
\begin{tabularx}{\linewidth}{@{}lrrL@{}}
\toprule
\textbf{Expense Item} & \textbf{Current (VND)} & \textbf{Target (VND)} & \textbf{Control Strategy} \\
\midrule
Fixed (Internet \& Phone)       & 320,000   & 320,000   & Maintain current plans. \\
Variable --- Food               & 3,810,866 & 3,000,000 & Reduce dining out and social drinking; increase home cooking. \\
Variable --- Sports (Badminton) & 759,000   & 1,000,000 & Health investment. Covers court bookings, shuttlecocks, restringing. \\
Variable --- Shopping online    & 404,999   & 200,000   & Restrict non-essential purchases. \\
Variable --- Entertainment      & 145,000   & 100,000   & Restrict gaming cafe visits and subscriptions. \\
Emergency Buffer                & 0         & 300,000   & Reserve for health or motorcycle maintenance. \\
\midrule
\textbf{Total Expenses} & \textbf{5,439,865} & \textbf{4,920,000} & \\
\bottomrule
\end{tabularx}
\end{table}

Note that sports is budgeted \textit{higher} than current spending (1,000,000 vs.\ 759,000~VND) to cover the full planned cost of court fees, shuttlecocks, and restringing each month. The budget follows a \textbf{35\% Needs / 15\% Wants / 50\% Savings} split. This directs about 4,580,000 VND each month---50.8\% of gross income---to savings, more than twice the standard 20\% benchmark. That margin is enabled almost entirely by not paying rent.

This housing advantage is the defining structural feature of the current budget. It will not last indefinitely: once I move out and pay my own rent---even a modest 5,000,000 to 7,000,000~VND/month room in Hanoi---the savings rate drops to 20 to 25\% at current income, which is still healthy but no longer exceptional. The period between now and when I pay my own rent is therefore the highest-leverage savings window available. Every month of this window that is used well compounds forward; every month wasted does not come back.

\section{Annual Budget Evaluation}

\begin{table}[H]
\caption{Annual Budget Evaluation (12-Month Projection)}
\label{tab:annual_budget}
\begin{tabularx}{\linewidth}{@{}lrrcL@{}}
\toprule
\textbf{Financial Metric} & \textbf{Planned (Year)} & \textbf{Projected (Current)} & \textbf{Met?} & \textbf{Variance Analysis} \\
\midrule
Income   & 114,000,000 & 114,000,000 & Yes & Stable internship, family, and freelance income assumed. \\
Expenses &  59,040,000 &  65,278,380 & No  & Food overspend is the entire source of the gap. \\
Savings  &  54,960,000 &  48,721,620 & Yes & High baseline; full target requires food discipline. \\
\bottomrule
\end{tabularx}
\end{table}

The 6,256,380~VND annual gap between current and target savings is entirely explained by excess food spending. No other category is off-target. Hitting the target budget produces \textbf{54,960,000~VND/year}---enough to fund the equipment upgrades (23,500,000~VND net) and begin the cash emergency fund within the same 12-month window.
